%1/2in left/right margins, 1in top, 1/2 bottom
\documentclass[12pt]{article}
\hbadness=10000
\tolerance=10000
\textheight 9.7in
\textwidth 7.0in
\oddsidemargin -.5in
\topmargin -.5in
\headsep 0in
\headheight 0in
%\pagestyle{empty}

\newcommand{\by}{\mbox{\boldmath$y$}}
\newcommand{\bX}{\mbox{\boldmath$X$}}
\newcommand{\bbeta}{\mbox{\boldmath$\beta$}}

\begin{document}
\hrule
\medskip
\centerline{\textbf{STATISTICS 801 - Section 250}}\hfill \textbf{Fall 2016}
\centerline{\textbf{Mini Project 4}}
\centerline{\textbf{Due Oct 10, 4:30 pm}}
\smallskip
\hrule
 Find a publication in an area that interests you, and presents a hypothesis test about a population mean, or compares 2 population means. Don't be surprised if the publication does not state all of the steps we discussed.
\begin{enumerate}
\item Print the publication, and provide the hard copy.
\item \textbf{Highlight} the problem they want to test. You have to be specific. Don't highlight a whole paragraph or page. Make a note next to the highlighted section that it is your hypothesis. 
\item State the null hypothesis in words.
\item State the null and alternative hypotheses using the formulas we learned in class.
\item \textbf{Highlight} every additional information (test statistics, alpha, p-value, decision, conclusion) that you can find in the publication.
\item State what steps of the hypothesis test are missing from the publication.
\item Try to fill in the gap. Sometimes there is not enough information given to calculate the test statistics, etc, but for example the conclusion will be stated so you can write the decision based on the conclusion. 
\item Justify why they used z or t. \textbf{Highlight} the evidence in the publication.
\end{enumerate}
Make sure you explain every highlighted section which problem they belong to.

\end{document}

